\section{Yukawa kernels and positivity lemmas}
\label{app:yukawa}

This appendix is the proof motor for the transport positivity, strong-CP protection, and bounded-residual statements used in the main text. The transport program is based on the regularized kernel
\[
\mathcal Y_f^{(\varepsilon)}=(D_f^\dagger D_f+\varepsilon^2)^{-1}.
\]
Three remarks are important.
\begin{enumerate}[leftmargin=1.5em, topsep=3pt, itemsep=3pt]
    \item The regulator $\varepsilon$ is not optional: without it the kernel becomes ill-defined exactly where the admissibility question is most delicate.
    \item Positivity of the kernel belongs to the admissibility-closure layer rather than to the hard carrier kernel. The main text closes it through the finite hexagon gap and the lift to the full admissibility-projected operator; this appendix records the transport-side mechanism behind that closure.
    \item Residual prefactors $\Lambda_{f,j}$ are bounded holonomy data inside the transport theorem. They must remain visible and must not be rhetorically hidden as if the mass map were diagonal without residual transport structure.
\end{enumerate}

\begin{remark}[Positive-kernel determinant route to strong CP]
If the physical mass map is written in the polar form
\[
M_f=U_{f,L}H_fU_{f,R}^\dagger,
\qquad
H_f>0,
\qquad
U_{f,L},U_{f,R}\in SU(3)_F,
\]
with $H_f$ induced by the positive transport kernel $\mathcal Y_f^{(\varepsilon)}$, then the induced determinant phase vanishes:
\[
\arg\det M_f=0.
\]
At the same time weak CP may survive through
\[
V_{\mathrm{CKM}}=U_{u,L}^\dagger U_{d,L}.
\]
Combined with the sheet-CP symmetry protection of $\theta_{\mathrm{sheet}}$ in the relative action, this is the algebraic core of the strong-CP theorem stated in the main text. The role of the appendix remark is to display the polar-decomposition mechanics behind that theorem, not to reopen its status.
\end{remark}

\section{Information-theoretic and cut-and-project readings}
\label{app:info-readings}

This appendix collects interpretive notes that help organize the algebraic structure of the carrier-first compiler. When an exact algebraic statement appears below, it refers back to a main-text proposition; the appendix itself does not upgrade any status label beyond the declared hard / closure / benchmark / conditional discipline.

\subsection{Integrity-index overdetermination of the canonical branch}

The canonical $41$-chain admits a compact integrity-table reading:

\begin{table}[H]
\centering
\small
\begin{tabularx}{\linewidth}{@{}>{\raggedright\arraybackslash}p{0.26\linewidth}>{\raggedleft\arraybackslash}p{0.12\linewidth}Y@{}}
\toprule
Layer & Value & Role \\
\midrule
Electroweak trace block $48\,I_1^{\mathrm{EW}}$ & $41$ & retained electroweak trace packet \\
Block exponent $32\,k_{\mathrm{EW}}$ & $41$ & E$_8$-side compression datum \\
UV hypercharge coefficient $10\,b_1$ & $41$ & one-loop electroweak closure datum \\
Occupancy-weighted carrier norm $\Oadm\gamma+1$ & $41$ & carrier / family compression identity \\
\bottomrule
\end{tabularx}
\caption{Integrity-index overdetermination of the canonical branch. The common integer is read as an algebraic integrity marker, not as a cryptographic statement.}
\label{tab:checksum-41}
\end{table}

The value of this table is structural rather than rhetorical: four independently defined layers collapse onto one integer on the retained branch. This is exactly the sense in which the branch carries the integrity index $\mathcal I_{41}$.

\subsection{Seed variance and surprisal readings}

The bridge identities also admit a disciplined probabilistic reading. Since
\[
\lambda_C^2=\phiz(1-\phiz),
\]
the Cabibbo scale is exactly the standard-deviation scale of a Bernoulli variable with success probability $p=\phiz$:
\[
\lambda_C=\sqrt{p(1-p)}.
\]
Equivalently, the Fisher information of the Bernoulli family is
\[
I(p)=\frac{1}{p(1-p)}=\frac{1}{\lambda_C^2},
\]
so the Cabibbo bridge measures the inverse square-root Fisher scale of the retained seed. This is the precise information-theoretic statement. The present paper does \emph{not} identify $\lambda_C$ itself with Fisher information or with a Bures distance.

Likewise, whenever the electromagnetic closure is written with
\[
\ln\!\bigl((\phiz)^{-1}\bigr),
\]
the exact probabilistic reading is the surprisal of the seed event. That term should not be called the Shannon entropy of the full Bernoulli distribution, which would instead be
\[
H(\phiz)
=
-\phiz\ln\phiz-(1-\phiz)\ln(1-\phiz).
\]
Finally, the neutrino bridge relation
\[
\sin^2\theta_{13}=\phiz e^{-\gamma}
\]
is naturally described as Boltzmann-like weighting by the carrier norm $\gamma$, but the present draft does not define a full thermodynamic partition function at that point. These readings are therefore organizational interpretations of already established formulas, not new closure theorems.

\subsection{Exact parity-code and information-compression reading}

The strongest information-theoretic statement is no longer a loose analogy. By \Cref{prop:parity-code-family}, the one-family packet is exactly the even-parity subspace on five occupation bits,
\[
S^+
\cong
\mathrm{span}\Bigl\{|x\rangle:\sum_{i=1}^5 x_i\equiv 0 \pmod 2\Bigr\},
\]
so the carrier packet is kinematically the classical even-parity $[5,4,2]$ code. The split formula
\[
S^+
=
\bigl(\Lambda^{\mathrm{even}}E_3\otimes\Lambda^{\mathrm{even}}E_2\bigr)
\oplus
\bigl(\Lambda^{\mathrm{odd}}E_3\otimes\Lambda^{\mathrm{odd}}E_2\bigr)
\]
shows that the coarse parity bits of the $3$-sector and $2$-sector are exactly locked. For the maximally mixed state on the $16$ basis words, the two parity observables are perfectly correlated and therefore carry exactly one bit of mutual information.

The split weight enumerator in \Cref{prop:split-weight-enumerator} is the corresponding code enumerator with hypercharge weights. It packages the six admissible occupancy sectors into the single polynomial
\[
W(u,v)
=
1+3u^2+6uv+v^2+2u^3v+3u^2v^2.
\]
The hypercharge-compression result of \Cref{prop:family-hypercharge-compression} sharpens the same picture further. On the uniform distribution over the $16$ carrier basis states, the six $Y$-sectors have probabilities
\[
\frac1{16}(1,3,6,1,2,3),
\]
so the hypercharge entropy is
\[
H(Y)
=
-\sum_i p_i\log_2 p_i
=
\frac72-\frac34\log_2 3
\approx
\num{2.311}\ \text{bits}.
\]
Since the total carrier word takes $4$ bits to specify, the unresolved nonabelian degeneracy is
\[
4-H(Y)\approx \num{1.689}\ \text{bits}.
\]
Thus hypercharge is an efficient family-sector compressor, but not a lossless label for the individual basis word.

\begin{remark}[Two code layers: kinematic and admissible]
The code reading should be split into two layers. First, $S^+$ is exactly the classical even-parity $[5,4,2]$ code on five occupation bits. Second, on the commuting admissible branch the protected physical sector is
\[
\mathcal H_{\mathrm{adm}}
:=
\operatorname{Ran}(\Padm)
=
\ker \Kadm
=
\ker \Delta_{\mathrm{adm}},
\]
where the last equality uses the Hodge description of \Cref{lem:hodge-admissibility-projector}. In that second layer the penalties $K_\Sigma,K_c,K_\Theta,K_g$ act as stabilizer-like sector constraints, and $\Padm$ is the Hodge projector onto the admissible physical subspace. This is the precise bridge to quantum-information language used in the manuscript. It is not yet a theorem of a self-correcting quantum memory, because no noise model, logical algebra, or lifetime analysis is derived here.
\end{remark}

\subsection{Conditional cut-and-project reading}

The split
\[
E=E_3\oplus E_2
\]
also admits a possible cut-and-project interpretation, with a three-dimensional visible sector and a two-dimensional internal sector. Such a reading would become mathematically sharp only after an explicit rank-$5$ lattice or module, a projection rule, and a seam/window selection theorem are written down. Until then, the quasicrystal language is only a conditional geometric reading of the carrier split rather than part of the hard core.

\section{Appendix elaboration on record algebra, horizons, and prediction semantics}
\label{app:record}

This appendix elaborates the main-text record and prediction constructions while keeping the more speculative horizon extrapolations outside the core theorem stack.

\subsection{Minimal admissible vacuum}

At finite volume or with a declared cutoff, an admissible vacuum candidate is written as
\[
\rho_B
:=
\argmin_{\rho\in\Vadm}\tr(\rho\Hrel),
\]
with lower boundedness and coercivity understood exactly as in the main-text closure axioms. The stronger generative theorem path described in the conclusion would refine this object to the global admissible vacuum functional $\mathfrak V[g,\chigeo,\Phi,\nabla_F,\delta,\{L_f\}]$ whose variations also select the transport branch. In this appendix, however, the same vacuum object is used only as the background state for the record-algebra elaborations; no stronger cosmological claim is added here beyond the closed existence statement already declared in the main theorem stack.

\subsection{Record growth principle}

One useful coarse-grained record algebra is
\[
\Arec(t):=\bigoplus_{\Delta_i>\Hhub(t)}\mathcal A_i,
\qquad
S_{\mathrm{rec}}(t)=\log\dim\Arec(t).
\]
As $\Hhub(t)$ decreases, more stable channels open. The intended arrow-of-time claim is therefore
\[
\frac{dS_{\mathrm{rec}}}{dt}>0,
\]
but only as a record-growth principle, not as a replacement for detailed nonequilibrium dynamics.

\subsection{Observer algebra and Born-type evaluation}

An observer is modeled as a stable pointer subalgebra
\[
\Aobs\subset\Arec.
\]
With orthogonal projectors $P_i\in\Aobs$ one writes the measurement channel
\[
\mathcal M_O(\rho)=\sum_i(P_i\rho P_i)\otimes |i\rangle\langle i|.
\]
In finite-dimensional sectors, positivity, additivity on orthogonal decompositions, and phase invariance lead to the standard Born-type assignment
\[
p_i=\tr(\rho P_i).
\]
The claim is therefore only a constrained evaluation principle on admissible sectors, not a new theorem about measurement in all settings.

\subsection{Horizons as conjectural inner seams}

If a horizon is modeled as an inner seam, the manuscript uses only an approximate factorization language,
\[
\mathcal H_{\mathrm{out}}\otimes\mathcal H_H\otimes\mathcal H_{\mathrm{in}},
\]
or, better, an algebraic split on an appropriate code subspace. This appendix intentionally avoids selling the factorization as exact fundamental structure.

\subsection{Prediction semantics}

The main-text prediction map can be decorated by explicit experimental metadata as
\[
Q=(\rho_0,\mathcal O_Q,\Delta_Q,t,\mathcal C_Q,\mathcal R_Q),
\qquad
\Pred(Q)=\mathcal R_Q\!\left(\tr(\rho_tE_Q(\Delta)),\mathcal C_Q\right),
\]
with
\[
\rho_t=U(t)\rho_0U^\dagger(t).
\]
This is kept only as an optional formal extension.

\section{Extended benchmarks}
\label{app:benchmarks}

This appendix turns the broader prediction surface into explicit experimental bookkeeping. The point is not to promote every row to main-text benchmark status, but to keep next test, kill criterion, and dependency structure visible wherever the theory already commits to a concrete target.

\subsection{Operational prediction matrix}

The appendix matrix keeps the same truth-mode discipline as the main text, but its final column adds a short row-level qualifier whenever a conditional branch, interface, or cross-check needs to be made explicit.

\begin{table}[H]
\centering
\tiny
\sisetup{round-mode=none, group-minimum-digits=20}
\setlength{\tabcolsep}{0.8pt}
\renewcommand{\arraystretch}{0.94}
\sloppy
\resizebox{\linewidth}{!}{%
\begin{tabular}{@{}>{\raggedright\arraybackslash}p{0.12\linewidth}>{\raggedright\arraybackslash}p{0.14\linewidth}>{\raggedright\arraybackslash}p{0.17\linewidth}>{\raggedright\arraybackslash}p{0.20\linewidth}>{\raggedright\arraybackslash}p{0.10\linewidth}>{\raggedright\arraybackslash}p{0.14\linewidth}@{}}
\toprule
\makecell[l]{Observable} & \makecell[l]{TFPT\\target} & \makecell[l]{Experimental reading\\next test} & \makecell[l]{Kill\\criterion} & \makecell[l]{Dependency\\class} & \makecell[l]{Status\\note} \\
\midrule
Cosmic birefringence $\beta$ & \num{0.2424}$^\circ$ & \makecell[l]{calibrated achromatic\\CMB rotation fit\\without EB self-nulling} & \makecell[l]{calibrated $\beta=0$\\within $\pm\num{0.05}^\circ$} & seed & \makecell[l]{Benchmark\\bridge} \\
$\delta_{\mathrm{CKM}}$ & \makecell[l]{\num{1.198}\\rad} & \makecell[l]{global flavor\\fit} & \makecell[l]{stable exclusion at\\$\ge 3\sigma$} & \makecell[l]{flavor\\bridge} & \makecell[l]{Benchmark\\bridge\\2nd order} \\
\makecell[l]{Rare kaons\\$K\to\pi\nu\bar\nu$} & \makecell[l]{$\mathrm{BR}(K^+)$\\$=\num{8.22e-11}$\\$\mathrm{BR}(K_L)$\\$=\num{2.65e-11}$} & \makecell[l]{NA62 / KOTO\\Grossman--Nir test} & \makecell[l]{$K^+$ outside\\$[6,10]\times10^{-11}$\\or $K_L$ mismatch} & \makecell[l]{flavor\\bridge} & \makecell[l]{Benchmark\\bridge} \\
\makecell[l]{PMNS phase /\\octant} & \makecell[l]{$\delta_{\mathrm{CP}}=240^\circ$\\$\sin^2\theta_{23}=\num{0.4557}$} & \makecell[l]{global oscillation\\fit} & \makecell[l]{exclude $240^\circ$\\or the lower octant\\at $\ge 3\sigma$} & \makecell[l]{neutrino\\branch} & \makecell[l]{Conditional\\extension\\oscillation\\branch} \\
$\Sigma m_\nu$ & \makecell[l]{\num{5.8764e-2}\\eV} & \makecell[l]{cosmological mass-sum\\analyses} & \makecell[l]{robust upper bound\\below the branch value} & \makecell[l]{neutrino\\branch} & \makecell[l]{Conditional\\extension\\oscillation\\branch} \\
$m_{\beta\beta}$ & \makecell[l]{\num{1.516e-3}\\eV} & \makecell[l]{light-Majorana\\$0\nu\beta\beta$ searches} & \makecell[l]{detection implying\\$m_{\beta\beta}\gtrsim \num{1e-2}$ eV} & \makecell[l]{neutrino\\branch} & \makecell[l]{Conditional\\extension\\Majorana\\interface} \\
\makecell[l]{nEDM / strong\\CP} & $\bar\theta=0$ & \makecell[l]{next-generation\\EDM searches} & \makecell[l]{stable nonzero\\hadronic EDM signal} & closure & Closure \\
\makecell[l]{Axion haloscope\\window} & \makecell[l]{$m_a\approx \num{65.19}$\\$\mu\mathrm{eV}$\\$\nu_a\approx \num{15.764}$\\$\mathrm{GHz}$} & \makecell[l]{haloscope scan at\\\num{15.764} GHz\\in the quoted window} & \makecell[l]{exclusion in\\\num{15.764} GHz\\$\pm \num{50}$ MHz} & \makecell[l]{scale\\grammar} & \makecell[l]{Conditional\\extension\\PQ interface} \\
$m_{\pi^0}$ & \makecell[l]{\num{134.979}\\MeV} & \makecell[l]{hadronic\\cross-check} & \makecell[l]{robust mismatch outside the\\stated uncertainty budget} & closure & \makecell[l]{Derived hadronic\\cross-check} \\
$\eta_B$ & \makecell[l]{\num{5.97e-10}} & \makecell[l]{optional leptogenesis\\branch} & \makecell[l]{robust exclusion of the\\quoted branch value} & \makecell[l]{scale\\grammar} & \makecell[l]{Conditional\\extension\\optional branch} \\
\bottomrule
\end{tabular}}
\caption{Operational prediction matrix extending the main-text prediction surface. Each row carries its own status note, experimental reading, kill criterion, and dependency class.}
\label{tab:extended-benchmarks}
\end{table}

\subsection{Target and kill-test rows}

\begin{table}[H]
\centering
\scriptsize
\sisetup{round-mode=none, group-minimum-digits=20}
\setlength{\tabcolsep}{2pt}
\renewcommand{\arraystretch}{0.96}
\resizebox{\linewidth}{!}{%
\begin{tabular}{@{}>{\raggedright\arraybackslash}p{0.22\linewidth}>{\raggedright\arraybackslash}p{0.24\linewidth}>{\raggedright\arraybackslash}p{0.12\linewidth}>{\raggedright\arraybackslash}p{0.38\linewidth}@{}}
\toprule
Test row & TFPT target or prohibition & Status reading & Failure mode \\
\midrule
Second light Higgs doublet & \makecell[l]{no additional seam-even\\light doublet} & \makecell[l]{kill\\test} & \makecell[l]{robust discovery would\\falsify the bosonic-index route} \\
Calibrated birefringence null & \makecell[l]{$\beta=\num{0.2424}^\circ$\\on the minimal branch} & \makecell[l]{kill\\test} & \makecell[l]{externally calibrated null\\consistent with zero at the\\stated tolerance would remove\\the minimal birefringence sector} \\
Rare-kaon corridor & \makecell[l]{$\mathrm{BR}(K^+)=\num{8.22e-11}$\\$\mathrm{BR}(K_L)=\num{2.65e-11}$} & \makecell[l]{kill\\test} & \makecell[l]{a stable kaon result outside\\the quoted corridor or\\incompatible GN-plane\\correlation would break the\\present flavor bridge} \\
Strong-CP / EDM signal & \makecell[l]{$\bar\theta=0$ on the\\admissible branch} & \makecell[l]{kill\\test} & \makecell[l]{stable nonzero EDM signal\\would break the present\\closure route} \\
Cosmological neutrino sum & \makecell[l]{$\Sigma m_\nu=\num{5.8764e-2}$ eV\\on the oscillation-inversion\\normal branch} & \makecell[l]{pressure\\test} & \makecell[l]{a robust upper bound below\\the stated branch value would\\force a neutrino-sector revision} \\
Haloscope window & \makecell[l]{$\nu_a\approx\num{15.76}$ GHz\\for the quoted branch target} & \makecell[l]{target\\row} & \makecell[l]{exclusion in the relevant\\coupled window would pressure\\the axion interface} \\
\bottomrule
\end{tabular}}
\caption{Target and kill-test rows. These are not current confirmations; they describe the most direct falsification or pressure points of the present branch.}
\label{tab:kill-tests}
\end{table}

\subsection{Axion interface}

The downstream axion rows should be read as an experimental interface rather than as an extra hard theorem. On the retained branch,
\[
f_a\approx \num{8.86e10}\,\mathrm{GeV},
\qquad
m_a\approx \num{65.19}\,\mu\mathrm{eV},
\qquad
\nu_a\approx \num{15.764}\,\mathrm{GHz}.
\]
In standard haloscope units this gives
\[
\left|g_{a\gamma\gamma}^{(\mathrm{phys})}\right|
=
\frac{|g_{a\gamma\gamma}|}{f_a}
\approx
\num{1.80e-12}\,\mathrm{GeV}^{-1},
\]
so the practical scan prescription is a conservative window
\[
\nu_a=\num{15.764}\,\mathrm{GHz}\pm \num{50}\,\mathrm{MHz}
\]
at the quoted coupled sensitivity. This interface depends on the downstream PQ / scale-grammar branch and should therefore not be read with the same independence status as the hard carrier kernel or the closure theorems.

\subsection{Supplementary mass-source rows}

The compact main-text mass ledgers intentionally suppress explicit formula detail. The source rows collected here retain that information without forcing the main text to read like a mixed formula notebook.

\begin{table}[H]
\centering
\scriptsize
\sisetup{round-mode=none, group-minimum-digits=20}
\setlength{\tabcolsep}{2pt}
\renewcommand{\arraystretch}{0.96}
\begin{tabularx}{\linewidth}{@{}>{\raggedright\arraybackslash}p{0.10\linewidth}>{\raggedright\arraybackslash}p{0.46\linewidth}>{\raggedleft\arraybackslash}p{0.16\linewidth}Y@{}}
\toprule
Particle & TFPT formula or source row & TFPT value & Status note \\
\midrule
$e$ & $\tfrac{\vgeo}{\sqrt2}\,\tfrac{16\pi}{7}(\phiz)^5$ & \makecell[r]{\num{0.52197}\\MeV} & \makecell[l]{constant-factory candidate\\uses $\vgeo$, not $\vewref$} \\
$\mu$ & $\tfrac{\vgeo}{\sqrt2}\,\tfrac{4\pi}{3}(\phiz)^3$ & \makecell[r]{\num{107.70}\\MeV} & \makecell[l]{constant-factory candidate\\uses $\vgeo$, not $\vewref$} \\
$\tau$ & $\tfrac{\vgeo}{\sqrt2}\,\tfrac{7\pi}{6}(\phiz)^2$ & \makecell[r]{\num{1.7723}\\GeV} & \makecell[l]{constant-factory candidate\\uses $\vgeo$, not $\vewref$} \\
\multicolumn{4}{@{}l}{\textit{Oscillation-inversion normal-ordering branch:}} \\
$\nu_1$ & oscillation-inversion normal-ordering branch & $0$ & retained branch value \\
$\nu_2$ & $\sqrt{\Delta m_{21}^2}$ & \makecell[r]{\num{8.614e-3}\\eV} & oscillation-inversion branch \\
$\nu_3$ & $\sqrt{\Delta m_{31}^2}$ & \makecell[r]{\num{5.015e-2}\\eV} & oscillation-inversion branch \\
$\Sigma m_\nu$ & \makecell[l]{$m_1+m_2+m_3$ on the\\oscillation-inversion branch} & \makecell[r]{\num{5.8764e-2}\\eV} & operational branch value used in Tables~11 and 20 \\
\multicolumn{4}{@{}l}{\textit{Geometric seesaw cross-check (separate from the operational branch):}} \\
$\nu_3^{\mathrm{ss}}$ & \makecell[l]{$e^{2\lambda_C^2}\vgeo^2/M_R$\\with $M_R\approx\num{1.31e15}$ GeV} & \makecell[r]{\num{4.939e-2}\\eV} & geometric seesaw cross-check \\
$\nu_2^{\mathrm{ss}}$ & $\pi\phiz\,m_3^{\mathrm{ss}}$ & \makecell[r]{\num{8.250e-3}\\eV} & geometric seesaw cross-check \\
$\Sigma m_\nu^{\mathrm{ss}}$ & $(1+\pi\phiz)\,m_3^{\mathrm{ss}}$ & \makecell[r]{\num{5.764e-2}\\eV} & separate from the operational branch sum \\
$u$ & ChPT ratio inversion anchored to $m_s$ & \makecell[r]{\num{2.195}\\MeV} & \makecell[l]{derived; $\overline{\mathrm{MS}}$\\at $2$ GeV} \\
$d$ & $m_s/(m_s/m_d)$ & \makecell[r]{\num{4.67}\\MeV} & \makecell[l]{derived; $\overline{\mathrm{MS}}$\\at $2$ GeV} \\
$s$ & strange-quark anchor in the light-quark solver & \makecell[r]{\num{93.4}\\MeV} & \makecell[l]{anchor-dependent;\\$\overline{\mathrm{MS}}$ at $2$ GeV} \\
$c$ & \makecell[l]{up-type ratio ladder with\\$m_c/m_u=\num{486.92}$} & \makecell[r]{\num{1.27}\\GeV} & threshold input fixes the absolute scale \\
$b$ & \makecell[l]{down-type ratio ladder with\\$m_b/m_s=\num{43.79}$} & \makecell[r]{\num{4.18}\\GeV} & threshold input fixes the absolute scale \\
$t$ & \makecell[l]{up-type ratio ladder with\\$m_t/m_c=\num{133.01}$} & \makecell[r]{\num{172.57}\\GeV} & \makecell[l]{pole / $\overline{\mathrm{MS}}$ distinction\\remains external} \\
$\pi^0$ & GMOR relation from $(m_u+m_d)$ and $\langle\bar qq\rangle$ & \makecell[r]{\num{134.979}\\MeV} & hadronic out-of-sample check \\
$p$ & hadron mass ladder & \makecell[r]{\num{935.03}\\MeV} & derived hadronic scale check \\
\bottomrule
\end{tabularx}
\caption{Supplementary mass-source rows referenced by the compact main-text mass surface.}
\label{tab:mass-source-supplement}
\end{table}

For the main-text baryon-fraction comparison we use the Planck 2018 base-$\Lambda$CDM conversion
\[
\Omega_b=\frac{\Omega_b h^2}{h^2},
\qquad
h=\frac{H_0}{100},
\]
so the quoted $\Omega_b$ row is not a free-floating cosmology number but an explicit reconstruction from the Planck parameter pair $(\Omega_b h^2,H_0)$ adopted in the suite ledger.

